L2的稀有符号事件说明，绝对误差还取决于函数的概率尺度。原参考与全部414份有效拟合均未访问$\theta_1>0$；另立协议的独立重要性参考给出两套估计$@L2_P1@$和$@L2_P2@$，各自近似MCSE及权重质量见补充材料第2节。对应全零估计的事件平方差仅约$10^{-18}$，相对于任一正概率参考的相对误差却为100\%。按有效拟合预算合计的\PBIIDLTwoStates{}个状态条目，即使假想全部独立且来自正确后验，以第一套参考计算也只期望约$\PBIIDLTwoHits$次命中，零命中概率约\PBIIDLTwoZero\%。实际链内相关且同核/跨设备执行共享轨迹，不能将此IID参照当作运行概率。全零本身不足以诊断混合不良；它说明普通事件频率在该预算下缺乏有用的相对精度，小绝对平方差掩盖了这一限制。新增参考单列，原18项失败、常量函数的未定诊断及全函数判据保持原样。
